\documentclass[11pt]{article}
\usepackage[a4paper,left=30mm,right=30mm,top=25mm,bottom=25mm,marginparwidth=16mm,marginparsep=4mm]{geometry}
\usepackage{fontspec,xcolor,needspace,longtable,array}
\usepackage[unicode,hidelinks]{hyperref}
\usepackage{polyglossia}
\setmainlanguage{english}
\setotherlanguages{arabic}
\setmainfont{Linux Libertine O}[Ligatures=TeX]
\newfontfamily\arabicfont{Amiri}[Script=Arabic]
\newfontfamily\ArNote{Amiri}[Script=Arabic,Scale=1.05]
\hypersetup{pdftitle={al-Battani, commentary on Ptolemy's Tetrabiblos: an English translation},pdfauthor={al-Battani}}
\setlength{\parindent}{0pt}\setlength{\parskip}{0.55em}
\tracinglostchars=2
\newcommand{\Fol}[1]{\marginpar{\footnotesize\color{gray}f.\,#1}}
\newcommand{\ChapterHead}[4]{\par\needspace{8\baselineskip}\bigskip\pdfbookmark[1]{#1}{ch-#2}\hypertarget{EN-K#2}{}%
 {\large\bfseries\interlinepenalty=10000 #1#4{\normalfont\normalsize#3}\par}\nobreak\medskip}
\begin{document}

\thispagestyle{empty}
\begin{center}\Large AL-BATTĀNĪ\par\vspace{4mm}\LARGE Commentary on Ptolemy's \emph{Tetrabiblos}\par\vspace{2mm}\large \emph{Sharḥ Kitāb al-Arbaʿ maqālāt fī aḥkām ʿilm al-nujūm}\par\vspace{5mm}\normalsize translated from the Arabic of the Escorial manuscript, árabe 969, ff. 20r–80v\end{center}\vspace{8mm}
\noindent This is an English translation of al-Battānī's commentary on the \emph{Tetrabiblos}, the earliest that survives. It is made from Nadine Löhr's transcription of the Escorial manuscript in \emph{Ptolemaeus Arabus et Latinus} (Bavarian Academy of Sciences and Humanities, 2021, CC BY 4.0). The commentary follows Ptolemy's work chapter by chapter, restating his doctrine and adding al-Battānī's own explanations, examples and procedures; the notes give the chapter of the \emph{Tetrabiblos} in F.~E. Robbins's numbering, which was used as a check on the meaning, and explain terms. Every sentence of the Arabic is translated, in the order of the text; numbers are given in figures. Words added for the sense stand in square brackets; words supplied by the editor of the transcription in angle brackets ⟨ ⟩. Words that the editor marks as additions of the copyist, such as invocations, stand in double brackets ⟦ ⟧, and words struck out by the copyist are left out. A doubtful reading of the transcription is marked (?). The folios of the manuscript are given in the margin.
\par\smallskip\noindent The translation is in progress. This version holds the title page and chapters 1, 2, 3, 4, 5, 6, 7, 8, 9. It is AI-integrated work (Claude), 2026; no person has proofread it. The translation is dedicated to the public domain (CC0 1.0); the Arabic it translates is Löhr's transcription, used under its licence (CC BY 4.0).
\par\clearpage
\begin{center}\Fol{20r}The Book of the Four Treatises on the Judgements of the Stars, by the lord of the sages, Ptolemy:\end{center}
\begin{center}the commentary of Abū ʿAbd Allāh Muḥammad ibn Jābir Abū Sinān al-Ḥarrānī, known as al-Battānī\footnote{So the manuscript, «Abū Sinān», for «ibn Sinān» (PAL: sic).}\end{center}
\begin{center}\Fol{20v}{\mdseries ⟦}In the name of God, the Merciful, the Compassionate; He is my sufficiency.{\mdseries ⟧}\end{center}
\ChapterHead{{\mdseries ⟨}1{\mdseries ⟩} On knowing the power of the sun, the moon, the wandering stars and the fixed stars}{01}{}{\footnote{Ptolemy, \emph{Tetrabiblos} I.4 and I.9. The chapters of the \emph{Tetrabiblos} are cited in F. E. Robbins's numbering.}}
\par Abū ʿAbd Allāh Muḥammad ibn Jābir al-Ḥarrānī, known as al-Battānī, said: The action of the substance of the sun is to dry and to heat, the drying being slight. As for the moon, most of the action of its power is moistening, and it changes bodies visibly, because it ripens and rots most of what is in them; it also shares a little in heat with the sun, because it receives light from it. Saturn acts mostly by cooling, and it dries a little, because of its cold and dryness. Jupiter is of a temperate mixture, for it heats and moistens together, and it produces fertilizing winds. Mars dries above all, and most of the action of its power is burning, because its nature is burning, its colour being like the colour of fire. Venus is of a temperate mixture, except that it moistens much and heats a little. As for Mercury, its moistening and its drying are equal: at some times it dries and absorbs moisture, at others it moistens, and it passes quickly from one state to the other; and it generates winds.\footnote{«Fire» reads \RL{\ArNote النار} with the Arabic \emph{Tetrabiblos} (PAL's T); the manuscript has \RL{\ArNote العار}, «shame».}\par
\par As for the powers of the fixed stars, from their mixtures, compared with the powers and natures of the wandering stars as they were found for them when they were observed: we have set them down in the tables of the positions of the fixed stars, in a book other than this, and drawn beside each star a mark by which the wandering stars that its nature resembles are known. Some resemble in their mixture the mixture of one wandering star; some resemble the mixture of two stars equally; and some resemble one of two stars much and the other a little. Their natures resemble [the wandering stars] as the marks found there show, and for that reason we can dispense with mentioning their powers here.\footnote{«Beside» reads \RL{\ArNote إزاء} for the transcription's \RL{\ArNote إذا}. Ptolemy gives the powers of the fixed stars in \emph{Tetrabiblos} I.9. The star catalogue of al-Battānī's \emph{Zīj}, as Nallino edits it from the Escorial codex, has no such marks.}\par
\ChapterHead{{\mdseries ⟨}2{\mdseries ⟩} On the powers of the planets according to the configurations they have relative to the sun}{02}{}{\footnote{\emph{Tetrabiblos} I.8.}}
\par \Fol{21r}The moon and the three upper planets receive increase and decrease in their particular powers according to the configurations they have with the sun. From the time when the moon appears as a crescent to the time of its first half in light, it produces moisture above all; from its first half to its full, when it is in opposition to the sun, heat; from its full to its second half in light, dryness; and from its second half to its disappearance, cold.\par
\par Saturn, Jupiter and Mars, from the time of their rising in the mornings and their appearing from under the rays of the sun to the time of their first station, moisten more; from their first station to their rising at the beginning of the night, which is when they are in opposition to the sun, they heat more; from their rising at the beginning of the night to their second station, they dry more; and from their second station to their setting and disappearance, they cool more. When they mix with one another, they produce in the air very many kinds of very different conditions, and these conditions are called the causes of the stars.\par
\ChapterHead{{\mdseries ⟨}3{\mdseries ⟩} On the powers of the four seasons of the year and of the ages of animals}{03}{}{\footnote{\emph{Tetrabiblos} I.10.}}
\par In spring moisture prevails, in summer heat, in autumn dryness and in winter cold. In the same way, in the first age of animals the moist substance prevails; in the second age, up to the time of the prime, heat; in the third age, when it nears the decline and begins to decrease, dryness; and in the last of the ages, when it nears dissolution, cold: like spring, summer, autumn and winter.\par
\ChapterHead{{\mdseries ⟨}4{\mdseries ⟩} On the powers of the four angles and the four quarters of the horizon}{04}{}{\footnote{\emph{Tetrabiblos} I.10.}}
\par So too is the condition of the four quarters of the horizon and of the angles from which the winds blow to all sides. Dryness prevails in the quarter next to the east, and the wind that blows from it, the east wind, which is called Apeliotes, is without moisture and drying. The winds that blow from the side of the south, which are called southern winds, are hot and burning. The quarter of the west is moist, and the winds that blow from it, which are called western, soften and moisten. The quarter of the north is very cold, and the winds that blow from this quarter, which are called northern, are cold and condensing.\par
\par \Fol{21v}The knowledge of these matters is also useful and helps to distinguish the mixtures. For it is clear that the actions of the power in the stars change with the conditions of the seasons of the year, of the ages or of the angles; that when the conditions are akin to them, their qualities are pure and unmixed, with nothing mingled in them, and their action is stronger, as the stars that heat by nature are stronger in hot conditions and the moist stars stronger in moist conditions; and that in the opposites of these conditions their power is not pure and unmixed but, being mixed with something else, weak, as the power of the heating stars is weaker in cold conditions and that of the moist stars weaker in dry conditions. The natures of the fixed stars that mix with those stars, or share with them in one of the nine configurations that they have, especially with the sun, which were mentioned earlier in this book, increase this, decrease it and change it not a little. In the same way the change in all the other conditions follows the mixing of the qualities.\par
\ChapterHead{{\mdseries ⟨}5{\mdseries ⟩} On the fixed, the solstitial and the bicorporeal signs}{05}{}{\footnote{\emph{Tetrabiblos} I.11.}}
\par Of the signs, two are solstitial and two are equinoctial. The solstitial ones are Cancer and Capricorn, and the equinoctial ones Aries and Libra. These signs are given these names because when the sun comes to the point of the summer solstice it brings about summer; when it comes to the point of the winter solstice, winter; when it comes to the point of the spring equinox, spring; and when it comes to the point of the autumn equinox, autumn; and at these two points night and day are equal over the whole earth.\par
\par Of the eight remaining signs, which lie between the points of the equinoxes and of the solstices, four are fixed. They are called fixed because the season stays fixed in them in one state; they are the ones that follow the two signs of the equinoxes and the two signs of the solstices, and they are called Taurus, Scorpio, Leo and Aquarius. The bicorporeal signs are those that follow the fixed signs. They are given this name because they lie between the fixed signs and the signs of the equinox and the solstice, and share, as it were, with their beginnings in the fixed signs and with their ends in the changing signs; they are Gemini, Virgo, Sagittarius and Pisces.\par
\ChapterHead{{\mdseries ⟨}6{\mdseries ⟩} On the sharing of the signs in the configurations}{06}{}{\footnote{\emph{Tetrabiblos} I.13–17.}}
\par \Fol{22r}As for the matters of the sharings in the configurations, which are the conjunction, the opposition, the quadrature, the sextile and the trine; of the signs that look at one another and are equal in power; of the commanding and the obeying signs; and of the signs that are not joined to one another and do not share in the configurations: we have explained them earlier in this book, since [there was given] the cause that shows why these distances alone are used and none other of the geometrical things. So too the others: the contrary and associated signs are those opposite one another; the hostile signs are those in quadrature to one another; and the friendly signs are those in trine to one another, and also those in sextile, since these two configurations fall between signs of one kind, either masculine or feminine, while the first two configurations are otherwise. Moreover, the houses of the two maleficent planets are in opposition and in quadrature to the houses of the luminaries, and the houses of the two beneficent planets in trine and in sextile to the houses of the luminaries.\footnote{Ptolemy explains why only these distances are used in \emph{Tetrabiblos} I.13, al-Battānī in chapter 54 of his \emph{Zīj}. «The two maleficent planets» reads \RL{\ArNote النحسين} for the transcription's \RL{\ArNote الجنسين}: the houses of Saturn, Capricorn and Aquarius, are opposite those of the moon and the sun, Cancer and Leo, and the houses of Mars, Aries and Scorpio, in quadrature to them (\emph{Tetrabiblos} I.17).}\par
\ChapterHead{{\mdseries ⟨}7{\mdseries ⟩} On the masculine and feminine signs, their other conditions and the names given to them}{07}{}{\footnote{\emph{Tetrabiblos} I.12. Ptolemy declines to list the epithets that al-Battānī gives below.}}
\par The signs feminine and masculine by nature we have explained in mentioning the triplicities. Some people have used another way of ordering the masculine and feminine signs, from the sign of the ascendant: they counted one masculine and one feminine, and made the beginning from the masculine; just as some people take the solstitial signs beginning from the sign of the moon, because the moon, through the speed of its motion, turns more than the other stars. Some of them divided the signs into four parts and held that the signs attributed to the mornings, the masculine ones, are those from the ascendant to the middle of the sky and those opposite them, from the setting point to the cardine of the earth; and they held that the two remaining quarters are attributed to the evenings and are feminine, so they made the signs that fall in these two quarters feminine and attributed them to the evenings.\par
\par \Fol{22v}They have given the signs many epithets, taken from the figures to which they are attributed: they called some of them four-footed, some wild, some governing, domestic, of much sowing, of little sowing and the like. They said that the domestic signs are Aries and Taurus; the wild, Leo and Sagittarius; the speaking signs, which resemble human beings, Gemini, Aquarius, Libra, Virgo and the first half of Sagittarius; the winged, Virgo, Sagittarius, Gemini and Pisces; the \emph{maʿqūma}, Aries, Taurus, Scorpio, Capricorn and Cancer; those of few offspring and little sowing, Aries, Taurus, Libra, Sagittarius, Capricorn and Aquarius; those of many offspring, Cancer, Scorpio and Pisces; the barren, Gemini, Leo and Virgo; those with mutilated limbs, Aries, Taurus, Leo and Scorpio; those of much anger, Aries and Leo; those that signify beauty and comeliness, Gemini, Libra, Virgo, Scorpio, Sagittarius and Pisces; those that gather and fill, Aries and its triplicity; those that give great wealth, Capricorn and its triplicity; those that pour out and empty, Libra and its triplicity; and those that grasp and take, Cancer and its triplicity. The signs of a strong voice are Gemini, Virgo and Libra; the moderate, which have half a voice, Aries, Taurus, Leo and Sagittarius; those of a weak voice, Capricorn and Aquarius; and those with no voice, Cancer, Pisces and Scorpio. Those that signify great care are Leo, Scorpio and Capricorn; those of little care, Virgo and Libra, for their darkness (?); and those of cunning and deceit, Leo, Sagittarius, Capricorn and Pisces. And they said that Aries, Taurus and Gemini, when their lords are found in them, {\mdseries ⟨}are{\mdseries ⟩} moving; Cancer, Leo and Virgo, when their lords are found [in them], are at rest; Libra, Scorpio and Sagittarius likewise, when their lords are found in them, are moving; and Capricorn, Aquarius and Pisces likewise, when their lords are found in them, are at rest.\footnote{The meaning of \RL{\ArNote المعقومة}, \emph{maʿqūma}, from the root of barrenness, is not clear; the barren signs follow under their usual name, \RL{\ArNote العقيمة}.}\par
\ChapterHead{{\mdseries ⟨}8{\mdseries ⟩} On the planets' becoming masculine and feminine relative to the sun, and also relative to the horizon}{08}{}{\footnote{\emph{Tetrabiblos} I.6.}}
\par \Fol{23r}They hold that the planets become masculine and feminine according to the configurations they have relative to the sun: when they are eastern, preceding the sun, they are masculine, and when they are western, following it, they are feminine. This also happens relative to the horizon: when they are in the positions from the east to the middle of the sky, or from the west to the cardine of the earth, they are masculine, because these two quarters are masculine; the first of them is called «the swift» and the second «the middle of speed». When they are between the middle of the sky and the west, or between the cardine of the earth and the ascendant, they are feminine, because these two quarters are feminine; the first of them is called «the middle of slowness» (?), and the second, between the ascendant and the cardine of the earth, «the slow and heavy».\footnote{The words from «they are masculine, because» to «or between the cardine of the earth» are written in the margin. «The swift» reads \RL{\ArNote السريع} for \RL{\ArNote الشربع}, and «the slow» \RL{\ArNote البطيء} for \RL{\ArNote البطن}, «the belly».}\par
\ChapterHead{{\mdseries ⟨}9{\mdseries ⟩} On facing, illumination and the like}{09}{}{\footnote{\emph{Tetrabiblos} I.23, where Robbins translates Ptolemy's terms «proper face» and «chariots».}}
\par It is said that the planets have a proper facing when each of them shares in a configuration with the sun and the moon in this way: that there is between it and the sun or the moon the same distance as between its house and the house of the sun or of the moon, except that it is western of the sun, I mean in the following signs, and eastern of the moon, in the preceding signs. For example, Venus in sextile to one of the luminaries, but in the third sign from the sign of the sun in the order of the signs, or in the eleventh sign from the sign in which the moon is; or Mercury in the second sign from the sign of the sun, or in the twelfth from the sign of the moon. For the other planets reckon their houses in the same way. When the planets are in this condition with respect to the luminaries, they are facing them.\footnote{The manuscript first wrote «or in the twelfth from the sign of the moon, and the other planets», struck it out, and wrote the correct passage in the margin, from «in the order of the signs» to «or in the twelfth».}\par
\par As for illumination, it is said [to be] when the planets are in their own domain and in places like it that agree with them; and when one of them has two or more of its dignities, as when it is in its house, its exaltation, its triplicity and its term and the like, [it is said] to be in the place of its joy and honour, and then its power increases in its actions that resemble the signs that belong to it and are akin to them in nature. It is also said that they rejoice when they are in the dignities of planets of their own domain; this is weaker than the first, because it is below it in power, but they share in it in this way all the same: just as, when they are found in foreign places or in a domain contrary to their own, their proper power weakens, and the contrariety of their mixture to the mixture of the signs in which they are becomes a cause for an action compounded of the two different natures to arise from them.\footnote{«Domain», \RL{\ArNote حيّز}, \emph{ḥayyiz}; Ptolemy speaks here of the «sect» of the planets.}\par
\end{document}
